\chapter[Business Architecture dan Tinjauan Pustaka v1]{Business Architecture: Memetakan Kapabilitas, Value Stream, dan Tinjauan Pustaka v1}
\label{ch:business-architecture}
\setlength{\emergencystretch}{2em}

\section*{Tujuan Pembelajaran}
\begin{itemize}
  \item Menjelaskan peran Business Architecture dalam menghubungkan strategi,
        operating model, dan kebutuhan perubahan arsitektur enterprise.
  \item Menyusun peta kapabilitas bisnis yang membedakan kapabilitas inti,
        pendukung, dan tata kelola.
  \item Memetakan value stream untuk menemukan titik masalah, kebutuhan data,
        aplikasi pendukung, dan peluang perbaikan.
  \item Mengubah artefak Business Architecture menjadi narasi analisis domain
        bisnis dan draf tinjauan pustaka v1 untuk makalah publikasi.
\end{itemize}

Tujuan pertama mendukung CLO 1 dan CLO 3. Tujuan kedua sampai keempat
mendukung CLO 3, CLO 6, dan CLO 7 karena mahasiswa diminta menyusun artefak
BDAT sekaligus menjadikannya argumen akademik dalam makalah individual.

\section{Ringkasan Awal Bab}

Bab ini membahas Business Architecture sebagai jembatan antara strategi
organisasi dan rancangan arsitektur yang lebih rinci. Pada Bab~\ref{ch:strategi-bisnis},
mahasiswa telah merumuskan masalah, pertanyaan riset, dan kontribusi awal.
Pada Bab~\ref{ch:togaf-lite}, mahasiswa menyusun visi arsitektur dan
bibliografi awal. Bab ini mengubah keduanya menjadi analisis domain bisnis:
kapabilitas apa yang harus dimiliki organisasi, alur nilai apa yang menjadi
sumber masalah, dan prioritas perubahan apa yang masuk akal.

Dua artefak utama yang digunakan adalah capability map dan value stream.
Capability map menunjukkan kemampuan organisasi yang perlu dipertahankan,
diperbaiki, atau dibangun. Value stream menunjukkan bagaimana nilai mengalir
dari kebutuhan pengguna sampai hasil yang diterima. Bagi pembaca non-TI,
kedua artefak ini penting karena membantu membahas arsitektur tanpa langsung
tenggelam dalam aplikasi, database, atau infrastruktur.

Keluaran bab ini adalah peta kapabilitas, value stream, prioritas kapabilitas,
dan draf Literature Review v1. Semua keluaran tersebut dipakai untuk
memperkuat bagian tinjauan pustaka, konteks studi kasus, dan analisis domain
bisnis dalam makalah publikasi individual.

\section{Pendahuluan}

Business Architecture menjawab pertanyaan yang tampak sederhana tetapi sering
terabaikan: kemampuan bisnis apa yang harus dimiliki organisasi agar strategi
digitalnya dapat berjalan? Pertanyaan ini berbeda dari pertanyaan tentang
aplikasi apa yang perlu dibangun. Aplikasi adalah salah satu pendukung.
Kemampuan bisnis mencakup proses, peran, aturan keputusan, data yang
dibutuhkan, mitra yang terlibat, ukuran keberhasilan, dan kendali risiko.

Dalam Enterprise Architecture, Business Architecture menjadi pintu masuk
untuk domain lain. Data Architecture membutuhkan kejelasan tentang informasi
apa yang diperlukan bisnis. Application Architecture membutuhkan pemahaman
tentang proses dan kapabilitas mana yang perlu didukung. Technology
Architecture membutuhkan konteks prioritas, risiko, dan skala layanan.
Karena itu, membahas data, aplikasi, dan teknologi sebelum memahami domain
bisnis sering menghasilkan rancangan yang tampak canggih tetapi tidak menjawab
masalah organisasi.

Untuk makalah publikasi, Business Architecture membantu mahasiswa membuat
analisis yang lebih meyakinkan. Penulis dapat menunjukkan bahwa masalah yang
dibahas tidak berdiri sendiri, melainkan muncul dari ketidakselarasan antara
strategi, kapabilitas, value stream, dan tata kelola. Rujukan tentang
Enterprise Architecture sebagai fondasi pelaksanaan strategi menekankan bahwa
arsitektur harus menata hubungan antara model operasi, proses, data, dan
keputusan bisnis \cite{ross2006enterprise}. Pandangan tentang organisasi
digital juga menekankan pentingnya kapabilitas yang dirancang sadar, bukan
hanya proyek teknologi yang berjalan terpisah \cite{ross2019designed}.

\section{Skenario Pembuka dan Pertanyaan Pemandu}

Bank Wanua telah memiliki visi arsitektur awal: memperkuat layanan digital
dengan integrasi data nasabah, orkestrasi proses produk, dan API yang
terkendali. Namun, begitu visi itu dibahas bersama unit bisnis, muncul
pertanyaan yang lebih praktis. Kapabilitas apa yang sebenarnya lemah? Apakah
masalah utama ada pada pemasaran digital, verifikasi nasabah, manajemen
produk, kepatuhan data, layanan mitra, atau analitik nasabah? Jika semua
dinyatakan penting, tidak ada prioritas yang dapat dipertanggungjawabkan.

Direktur bisnis meminta peta sederhana: kemampuan apa yang sudah cukup baik,
kemampuan apa yang perlu diperbaiki, dan kemampuan apa yang harus dibangun
agar Bank Wanua dapat menjadi mitra digital UMKM lokal. Unit kepatuhan
menambahkan pertanyaan lain: pada alur layanan mana data pribadi paling banyak
berpindah tangan? Tim teknologi juga meminta kejelasan: proses mana yang
sebenarnya perlu didukung aplikasi baru, dan proses mana yang dapat diperbaiki
dengan perubahan prosedur atau tata kelola?

Di titik ini, diskusi tidak dapat lagi berhenti pada slogan strategi. Bank
Wanua membutuhkan Business Architecture yang dapat dibaca bersama oleh
pimpinan, unit bisnis, kepatuhan, dan tim teknologi. Mahasiswa yang menulis
makalah tentang kasus ini juga membutuhkan hal yang sama: artefak yang
mengubah masalah organisasi menjadi analisis yang runtut.

\subsection*{Pertanyaan Pemandu}

\begin{enumerate}
  \item Kapabilitas bisnis apa yang paling menentukan keberhasilan visi
        arsitektur?
  \item Bagaimana capability map membantu membedakan masalah inti dari gejala?
  \item Value stream mana yang paling relevan dengan pertanyaan riset?
  \item Bukti apa yang diperlukan agar analisis domain bisnis tidak sekadar
        berdasarkan dugaan?
  \item Bagaimana hasil Business Architecture masuk ke tinjauan pustaka dan
        analisis makalah publikasi?
\end{enumerate}

\section{Mengapa Bab Ini Penting bagi Pembaca Non-TI}

Pembaca non-TI tidak harus menjadi perancang sistem teknis untuk dapat
berkontribusi pada arsitektur enterprise. Justru, banyak keputusan penting
berada di wilayah bisnis: siapa pengguna prioritas, nilai apa yang dijanjikan,
proses apa yang harus distandardisasi, data apa yang harus dipercaya,
pengendalian risiko seperti apa yang wajib ada, dan kapabilitas mana yang
layak didanai terlebih dahulu.

Business Architecture memberi bahasa bersama untuk diskusi tersebut. Dengan
capability map, mahasiswa dapat membicarakan kemampuan organisasi tanpa
terjebak pada nama aplikasi. Dengan value stream, mahasiswa dapat menjelaskan
di bagian mana nilai tersendat, risiko muncul, atau pengalaman pengguna
menurun. Dengan dua artefak ini, rekomendasi dalam makalah menjadi lebih
mudah dipertanggungjawabkan karena pembaca dapat melihat hubungan antara
masalah, bukti, prioritas, dan implikasi.

\section{Konsep Inti dan Glosarium Mini}

\subsection{Business Architecture}

Business Architecture adalah deskripsi terstruktur tentang bagaimana organisasi
menciptakan nilai melalui kapabilitas, proses, informasi, peran, kebijakan,
dan hubungan dengan pemangku kepentingan. Dalam TOGAF, Business Architecture
menjadi salah satu domain utama yang membantu menerjemahkan visi arsitektur ke
dalam struktur bisnis yang dapat dianalisis \cite{togaf10}.

Secara manajerial, Business Architecture dapat dibayangkan sebagai peta
kemampuan organisasi. Peta ini tidak menggambarkan struktur jabatan, tetapi
menggambarkan apa yang harus mampu dilakukan organisasi. Dua unit dapat
berbeda dalam struktur organisasi, tetapi membutuhkan kapabilitas yang sama.
Sebaliknya, satu unit dapat memegang beberapa kapabilitas sekaligus. Karena
itu, capability map biasanya lebih stabil daripada bagan organisasi.

\subsection{Business Capability}

Business capability adalah kemampuan yang dibutuhkan organisasi untuk
menciptakan nilai atau menjalankan mandatnya. Contoh kapabilitas pada bank
daerah adalah akuisisi nasabah UMKM, penilaian kelayakan kredit, manajemen
risiko, pengelolaan data nasabah, layanan kanal digital, pengelolaan mitra,
dan kepatuhan regulasi.

Kapabilitas tidak sama dengan proses. Proses menjelaskan urutan aktivitas.
Kapabilitas menjelaskan kemampuan yang harus ada agar aktivitas dapat berjalan
dengan baik. Satu kapabilitas dapat muncul dalam beberapa proses, dan satu
proses dapat memerlukan beberapa kapabilitas. Perbedaan ini penting karena
makalah arsitektur tidak cukup hanya menggambar alur kerja; ia perlu
menjelaskan kemampuan organisasi yang perlu berubah.

\subsection{Capability Map}

Capability map adalah visualisasi kapabilitas organisasi. Peta ini biasanya
dikelompokkan ke dalam lapisan seperti kapabilitas strategis, kapabilitas inti
pelayanan, dan kapabilitas pendukung. Untuk makalah individual, capability map
tidak harus lengkap seperti dokumen konsultan. Yang lebih penting adalah peta
itu relevan dengan pertanyaan riset dan cukup jelas untuk menjelaskan prioritas
analisis.

\begin{figure}[htbp]
\centering
\scalebox{0.86}{
\begin{tikzpicture}[
  capnode/.style={rectangle, rounded corners=1.5mm, draw=praditagreen,
    fill=praditagreen!10, thick, align=center, text width=30mm,
    minimum height=10mm},
  core/.style={rectangle, rounded corners=1.5mm, draw=orange!80!black,
    fill=orange!18, thick, align=center, text width=30mm,
    minimum height=10mm},
  support/.style={rectangle, rounded corners=1.5mm, draw=blue!70!black,
    fill=blue!10, thick, align=center, text width=30mm,
    minimum height=10mm},
  layer/.style={font=\small\bfseries, align=right}
]
  \node[layer] at (-1.0,2.6) {Strategis\\dan tata kelola};
  \node[capnode] (s1) at (1.0,2.6) {Arah produk\\digital};
  \node[capnode] (s2) at (4.4,2.6) {Manajemen\\risiko};
  \node[capnode] (s3) at (7.8,2.6) {Tata kelola\\data};

  \node[layer] at (-1.0,1.0) {Kapabilitas\\inti};
  \node[core] (c1) at (1.0,1.0) {Akuisisi\\UMKM};
  \node[core] (c2) at (4.4,1.0) {Penilaian\\kebutuhan};
  \node[core] (c3) at (7.8,1.0) {Layanan\\produk};

  \node[layer] at (-1.0,-0.6) {Kapabilitas\\pendukung};
  \node[support] (p1) at (1.0,-0.6) {Operasi\\cabang};
  \node[support] (p2) at (4.4,-0.6) {Kemitraan\\digital};
  \node[support] (p3) at (7.8,-0.6) {Analitik\\nasabah};

  \draw[gray!45, rounded corners=2mm] (-2.6,-1.3) rectangle (9.4,3.25);
\end{tikzpicture}
}
\caption{Contoh sederhana capability map untuk layanan digital Bank Wanua.}
\label{fig:capability-map-bank-wanua}
\end{figure}

\subsection{Value Stream}

Value stream adalah alur penciptaan nilai dari kebutuhan awal sampai hasil
yang diterima pengguna atau pemangku kepentingan. Jika capability map
menunjukkan kemampuan yang dibutuhkan, value stream menunjukkan bagaimana
kemampuan itu bekerja bersama dalam satu perjalanan nilai.

Pada Bank Wanua, value stream untuk layanan UMKM dapat dimulai dari menarik
minat calon nasabah, memahami kebutuhan, dan menilai kelayakan. Tahap
berikutnya adalah menawarkan produk, melakukan persetujuan, mengaktifkan
layanan, lalu memantau dan mendukung nasabah. Di setiap tahap, penulis dapat
menanyakan: siapa aktornya, data apa yang digunakan, sistem apa yang
mendukung, risiko apa yang muncul, dan metrik apa yang menunjukkan
keberhasilan.

\subsection{Operating Model dalam Business Architecture}

Operating model menjelaskan tingkat integrasi data dan standardisasi proses
yang dibutuhkan organisasi untuk menjalankan strategi \cite{ross2006enterprise}.
Bab~\ref{ch:strategi-bisnis} telah memperkenalkan konsep ini sebagai alat
membaca strategi. Dalam Bab 4, operating model dipakai lebih operasional:
untuk menilai apakah kapabilitas tertentu perlu distandardisasi, diintegrasikan,
atau dibiarkan berbeda sesuai konteks unit.

Bank Wanua, misalnya, mungkin perlu data nasabah UMKM yang terintegrasi lintas
cabang, tetapi proses pendampingan UMKM boleh tetap menyesuaikan karakter
wilayah. Keputusan semacam ini adalah keputusan Business Architecture. Ia
menentukan apa yang harus seragam, apa yang boleh lokal, dan apa yang harus
terhubung melalui tata kelola data atau API.

\subsection*{Glosarium Mini}

\begin{description}
  \item[Business Architecture] Deskripsi terstruktur tentang kapabilitas,
        proses, informasi, peran, dan kebijakan yang memungkinkan organisasi
        menciptakan nilai.
  \item[Business capability] Kemampuan organisasi untuk menjalankan fungsi
        tertentu, misalnya akuisisi nasabah, manajemen risiko, atau tata
        kelola data.
  \item[Capability map] Peta kapabilitas yang menunjukkan kemampuan inti,
        pendukung, dan tata kelola yang relevan dengan strategi.
  \item[Value stream] Alur penciptaan nilai dari kebutuhan pemangku kepentingan
        sampai hasil yang diterima.
  \item[Pain point] Titik dalam alur nilai yang menimbulkan kelambatan, risiko,
        biaya, duplikasi, atau pengalaman buruk.
  \item[Prioritas kapabilitas] Keputusan tentang kapabilitas mana yang perlu
        dipertahankan, diperbaiki, atau dibangun terlebih dahulu.
\end{description}

\section{Kerangka Berpikir}

Kerangka Bab 4 bergerak dari visi arsitektur menuju analisis domain bisnis.
Pertanyaan riset dan visi arsitektur memberi arah. Capability map menunjukkan
kemampuan yang dibutuhkan. Value stream membantu menemukan titik masalah.
Analisis prioritas menentukan kapabilitas mana yang paling penting. Hasilnya
menjadi bahan analisis domain bisnis dan draf tinjauan pustaka v1.

\begin{figure}[htbp]
\centering
\scalebox{0.86}{
\begin{tikzpicture}[
  box/.style={rectangle, rounded corners=1.5mm, draw=praditagreen,
    fill=green!8, thick, align=center, text width=31mm, minimum height=11mm},
  outnode/.style={rectangle, rounded corners=1.5mm, draw=orange!80!black,
    fill=orange!16, thick, align=center, text width=33mm, minimum height=11mm},
  link/.style={-Latex, thick, draw=praditagreen}
]
  \node[box] (vision) {Visi arsitektur\\dan RQ};
  \node[box, right=7mm of vision] (cap) {Capability\\map};
  \node[box, right=7mm of cap] (stream) {Value\\stream};
  \node[box, below=8mm of stream] (priority) {Prioritas\\kapabilitas};
  \node[outnode, left=7mm of priority] (analysis) {Analisis domain\\bisnis};
  \node[outnode, left=7mm of analysis] (lit) {Literature\\Review v1};

  \draw[link] (vision) -- (cap);
  \draw[link] (cap) -- (stream);
  \draw[link] (stream) -- (priority);
  \draw[link] (priority) -- (analysis);
  \draw[link] (analysis) -- (lit);
  \draw[link, dashed] (lit) -- (vision);
\end{tikzpicture}
}
\caption{Alur berpikir Business Architecture untuk makalah publikasi.}
\label{fig:business-architecture-flow}
\end{figure}

Alur ini tidak harus selalu linear. Saat menyusun capability map, mahasiswa
dapat menemukan bahwa pertanyaan risetnya terlalu teknis. Saat menggambar
value stream, mahasiswa dapat menyadari bahwa masalah utama bukan aplikasi,
melainkan aturan keputusan yang tidak jelas. Saat menulis Literature Review
v1, mahasiswa dapat menemukan konsep yang membuat prioritas kapabilitas perlu
diubah. Revisi semacam ini adalah bagian wajar dari penulisan akademik.

\section{Konteks Indonesia}

Business Architecture di Indonesia perlu memperhatikan variasi kematangan
organisasi. Banyak organisasi memiliki strategi digital yang ambisius, tetapi
kapabilitas bisnisnya belum merata. Cabang atau unit daerah dapat memiliki
cara kerja berbeda. Data pelanggan atau warga sering tersimpan di beberapa
sistem. Proses persetujuan masih bergantung pada dokumen manual. Peran vendor
besar, sementara kemampuan internal untuk mengelola perubahan belum selalu
kuat.

Regulasi juga memengaruhi prioritas kapabilitas. Organisasi yang mengelola
data pribadi perlu memiliki kapabilitas pengelolaan persetujuan, pembatasan
akses, pencatatan penggunaan data, dan penanganan insiden yang sejalan dengan
perlindungan data pribadi \cite{uupdp2022}. Penyelenggara sistem elektronik
perlu memperhatikan kewajiban pengelolaan sistem dan data elektronik
\cite{pp712019}. Organisasi sektor publik yang bertukar data perlu memahami
prinsip Satu Data Indonesia \cite{satudata}. Lembaga keuangan yang membuka
integrasi pembayaran atau kemitraan digital perlu memperhatikan standar Open
API pembayaran nasional \cite{snapbi}.

Konteks lokal juga memengaruhi value stream. Di banyak sektor, pengalaman
pengguna tidak hanya terjadi di kanal digital. Nasabah bank daerah, pasien
rumah sakit, petani, mahasiswa, atau pelaku UMKM masih berinteraksi melalui
kantor cabang, petugas lapangan, agen, WhatsApp, dan dokumen fisik. Karena
itu, Business Architecture yang baik tidak memaksakan narasi serba aplikasi.
Ia membaca alur nilai secara utuh: digital, fisik, manusia, mitra, dan aturan.

\section{Studi Kasus Utama: Business Architecture Bank Wanua}

Visi Bank Wanua adalah menjadi mitra keuangan digital bagi UMKM lokal. Agar
visi tersebut tidak berhenti sebagai slogan, mahasiswa perlu memetakan
kapabilitas yang membuat visi itu mungkin. Peta awal dapat disusun dari tiga
pertanyaan:

\begin{enumerate}
  \item Kapabilitas apa yang langsung berhubungan dengan nilai bagi UMKM?
  \item Kapabilitas apa yang menjaga risiko, kepatuhan, dan keberlanjutan?
  \item Kapabilitas apa yang mendukung integrasi digital dan pembelajaran
        organisasi?
\end{enumerate}

\subsection{Capability Map Bank Wanua}

\begin{table}[htbp]
\centering
\scriptsize
\caption{Capability map awal Bank Wanua untuk layanan digital UMKM.}
\label{tab:capability-bank-wanua}
\begin{tabularx}{\textwidth}{@{}p{0.22\textwidth} p{0.28\textwidth} X X@{}}
\hline
\textbf{Kelompok} & \textbf{Kapabilitas} & \textbf{Kondisi awal} & \textbf{Prioritas} \\
\hline
Strategis dan tata kelola & Perencanaan produk digital & Produk dirancang per
inisiatif, belum memakai portofolio layanan yang jelas & Diperbaiki \\
Strategis dan tata kelola & Manajemen risiko dan kepatuhan & Kontrol kuat,
tetapi sering masuk terlambat dalam rancangan layanan & Diperbaiki \\
Strategis dan tata kelola & Tata kelola data nasabah & Definisi data dan
kepemilikan belum konsisten lintas kanal & Dibangun \\
Kapabilitas inti & Akuisisi nasabah UMKM & Masih bergantung pada cabang dan
relasi lokal & Dipertahankan dan diperkuat \\
Kapabilitas inti & Penilaian kebutuhan dan kelayakan & Data tersebar; asesmen
sering berulang & Diperbaiki \\
Kapabilitas inti & Layanan produk dan transaksi & Kanal digital tersedia,
tetapi integrasi proses belum mulus & Diperbaiki \\
Kapabilitas pendukung & Kemitraan digital & Peluang mitra ada, tetapi aturan
integrasi belum matang & Dibangun \\
Kapabilitas pendukung & Analitik nasabah & Laporan tersedia, tetapi belum
menjadi dasar personalisasi layanan & Dibangun \\
\hline
\end{tabularx}
\end{table}

Tabel~\ref{tab:capability-bank-wanua} tidak dimaksudkan sebagai daftar final.
Ia adalah peta kerja awal. Dalam makalah, mahasiswa perlu menjelaskan dari
mana penilaian kondisi awal diperoleh: dokumen publik, wawancara, observasi,
laporan tahunan, studi literatur, atau asumsi yang dinyatakan secara eksplisit
bila kasus bersifat hipotetis-realistis.

\subsection{Value Stream Layanan UMKM}

Setelah kapabilitas dipetakan, mahasiswa perlu melihat bagaimana kapabilitas
tersebut bekerja dalam alur nilai. Untuk Bank Wanua, salah satu value stream
yang relevan adalah ``UMKM memperoleh layanan pembiayaan dan pembayaran
digital''. Alur ini tidak hanya mencakup aplikasi mobile, tetapi juga
interaksi cabang, verifikasi dokumen, penilaian risiko, aktivasi layanan, dan
pendampingan.

\begin{figure}[htbp]
\centering
\scalebox{0.78}{
\begin{tikzpicture}[
  stage/.style={rectangle, rounded corners=1.5mm, draw=orange!80!black,
    fill=orange!16, thick, align=center, text width=25mm, minimum height=12mm},
  pain/.style={rectangle, rounded corners=1.5mm, draw=red!70!black,
    fill=red!8, align=center, text width=25mm, minimum height=9mm,
    font=\footnotesize},
  link/.style={-Latex, thick, draw=orange!80!black}
]
  \node[stage] (s1) {Menarik\\minat UMKM};
  \node[stage, right=5mm of s1] (s2) {Memahami\\kebutuhan};
  \node[stage, right=5mm of s2] (s3) {Menilai\\kelayakan};
  \node[stage, below=11mm of s3] (s4) {Menyetujui\\produk};
  \node[stage, left=5mm of s4] (s5) {Mengaktifkan\\layanan};
  \node[stage, left=5mm of s5] (s6) {Memantau dan\\mendukung};

  \draw[link] (s1) -- (s2);
  \draw[link] (s2) -- (s3);
  \draw[link] (s3) -- (s4);
  \draw[link] (s4) -- (s5);
  \draw[link] (s5) -- (s6);

  \node[pain, below=4mm of s2] {Data kebutuhan\\belum terpadu};
  \node[pain, above=4mm of s3] {Asesmen\\berulang};
  \node[pain, below=4mm of s5] {Aktivasi kanal\\belum mulus};
\end{tikzpicture}
}
\caption{Value stream awal layanan UMKM Bank Wanua.}
\label{fig:value-stream-bank-wanua}
\end{figure}

Gambar~\ref{fig:value-stream-bank-wanua} membantu mahasiswa menemukan hubungan
antara kapabilitas dan titik masalah. Jika data kebutuhan UMKM belum terpadu,
kapabilitas tata kelola data dan analitik nasabah perlu diperkuat. Jika
asesmen berulang, kapabilitas penilaian kelayakan dan orkestrasi proses perlu
diperbaiki. Jika aktivasi kanal belum mulus, kapabilitas layanan produk,
integrasi aplikasi, dan tata kelola API mulai terlihat sebagai kebutuhan bab
berikutnya.

\section{Langkah Analisis dan Artefak Pendukung}

Bagian ini menjadi panduan kerja individual untuk menghasilkan artefak Bab 4.
Gunakan organisasi studi kasus yang sama dengan bab sebelumnya.

\subsection{Langkah 1: Tetapkan Ruang Lingkup Domain Bisnis}

Mulailah dari pertanyaan riset dan visi arsitektur. Tentukan layanan, produk,
unit, atau proses yang akan dianalisis. Jangan mencoba memetakan seluruh
organisasi jika makalah hanya membahas satu masalah. Ruang lingkup yang baik
cukup sempit untuk dianalisis, tetapi cukup penting untuk menjawab RQ.

\subsection{Langkah 2: Daftarkan Kapabilitas Bisnis}

Tuliskan kapabilitas yang dibutuhkan agar strategi berjalan. Gunakan kata
benda abstrak yang menunjukkan kemampuan, bukan aktivitas terlalu rinci.
Misalnya, tulis ``manajemen risiko kredit UMKM'', bukan ``mengisi formulir
penilaian''. Tulis ``tata kelola data pelanggan'', bukan ``mengunduh file
Excel''. Kapabilitas yang baik membantu pembaca melihat kemampuan organisasi
yang perlu ditata.

\subsection{Langkah 3: Kelompokkan Kapabilitas}

Kelompokkan kapabilitas ke dalam tiga sampai empat kategori. Kategori umum
adalah strategis dan tata kelola, kapabilitas inti, kapabilitas pendukung, dan
kapabilitas pembelajaran organisasi. Pengelompokan ini membantu pembaca non-TI
memahami mana kemampuan yang langsung menciptakan nilai dan mana yang menjaga
agar nilai itu dapat diberikan secara aman dan berulang.

\subsection{Langkah 4: Nilai Kondisi Awal dan Prioritas}

Untuk setiap kapabilitas, beri penilaian kondisi awal. Gunakan kategori
sederhana: kuat, cukup, lemah, atau belum ada. Setelah itu, tuliskan prioritas:
dipertahankan, diperbaiki, atau dibangun. Hindari memberi semua kapabilitas
prioritas tinggi. Prioritas yang baik harus berani memilih.

\subsection{Langkah 5: Pilih Value Stream Utama}

Pilih satu value stream yang paling relevan dengan RQ. Value stream dapat
berupa perjalanan pengguna, alur layanan, alur persetujuan, atau alur
penciptaan nilai. Tuliskan tahap utama dari awal sampai akhir. Pada setiap
tahap, catat aktor, data, aplikasi, risiko, dan pain point.

\subsection{Langkah 6: Hubungkan Value Stream dengan Kapabilitas}

Tanyakan kapabilitas mana yang mendukung setiap tahap value stream. Hubungan
ini penting karena menunjukkan mengapa sebuah kapabilitas perlu diprioritaskan.
Jika satu kapabilitas muncul di banyak titik masalah, kapabilitas itu mungkin
lebih penting daripada yang awalnya terlihat.

\subsection{Langkah 7: Susun Bukti dan Literatur Pendukung}

Business Architecture tidak boleh hanya berdasarkan intuisi. Catat sumber
bukti untuk setiap klaim. Bukti dapat berupa dokumen organisasi, laporan
tahunan, wawancara, observasi, regulasi, standar, atau literatur akademik.
Gunakan literatur untuk menjelaskan konsep seperti Enterprise Architecture,
operating model, kapabilitas digital, platform, dan tata kelola.

\subsection{Langkah 8: Tulis Narasi Analisis Domain Bisnis}

Ubah capability map dan value stream menjadi narasi. Narasi minimal menjawab:
kapabilitas apa yang paling bermasalah, bagaimana masalah itu terlihat dalam
alur nilai, mengapa masalah itu penting bagi strategi, dan apa implikasinya
untuk domain data, aplikasi, teknologi, API, dan platform pada bab berikutnya.

\begin{table}[htbp]
\centering
\scriptsize
\caption{Artefak Bab 4 dan pemakaiannya dalam makalah.}
\label{tab:artefak-bab4}
\begin{tabularx}{\textwidth}{@{}p{0.24\textwidth} p{0.32\textwidth} X@{}}
\hline
\textbf{Artefak} & \textbf{Isi minimum} & \textbf{Dipakai untuk bagian makalah} \\
\hline
Capability map & Kelompok kapabilitas, kondisi awal, prioritas & Konteks
studi kasus dan analisis domain bisnis. \\
Value stream & Tahap alur nilai, aktor, data, pain point, risiko & Analisis
baseline dan pembahasan masalah. \\
Tabel prioritas kapabilitas & Kapabilitas, alasan prioritas, bukti, dampak &
Gap analysis dan implikasi manajerial. \\
Matriks pustaka bisnis & Konsep, sumber, fungsi dalam argumen, hubungan ke RQ
& Literature Review v1. \\
Narasi analisis bisnis & Dua sampai empat paragraf akademik & Bagian
analysis/results atau studi kasus. \\
\hline
\end{tabularx}
\end{table}

\section{Integrasi ke Makalah Publikasi}

Bab ini memperkuat dua bagian utama makalah: tinjauan pustaka dan analisis
domain bisnis. Pada tinjauan pustaka, mahasiswa menjelaskan konsep Enterprise
Architecture, Business Architecture, operating model, kapabilitas digital, dan
value stream. Pada analisis domain bisnis, mahasiswa menunjukkan bagaimana
konsep tersebut dipakai untuk membaca kasus.

\subsection{Dari Artefak ke Tinjauan Pustaka}

Literature Review v1 tidak perlu langsung sempurna. Pada tahap ini, tujuannya
adalah membuat sintesis awal yang terhubung dengan artefak. Misalnya, jika
capability map menunjukkan lemahnya tata kelola data dan analitik nasabah,
tinjauan pustaka perlu memuat konsep tentang operating model, kapabilitas
digital, dan arsitektur enterprise sebagai fondasi eksekusi strategi
\cite{ross2006enterprise,ross2019designed}. Jika value stream menunjukkan
kebutuhan kemitraan dan layanan API, tinjauan pustaka mulai dapat menyinggung
literatur API dan platform sebagai pengantar untuk bab berikutnya
\cite{jacobson2011apis,parker2016platform}.

\subsection{Struktur Literature Review v1}

Untuk Bab 4, Literature Review v1 dapat disusun dengan pola berikut:

\begin{enumerate}
  \item paragraf tentang Enterprise Architecture sebagai penghubung strategi
        dan eksekusi;
  \item paragraf tentang Business Architecture, capability map, dan value
        stream;
  \item paragraf tentang konteks digital organisasi, misalnya platform, API,
        atau tata kelola data sesuai topik;
  \item paragraf sintesis yang menjelaskan celah atau kebutuhan analisis pada
        studi kasus mahasiswa.
\end{enumerate}

Contoh paragraf sintesis:

\begin{quote}
Literatur Enterprise Architecture menekankan bahwa strategi digital perlu
diterjemahkan ke dalam kapabilitas, proses, dan model operasi yang konsisten.
Dalam kasus bank daerah, kebutuhan integrasi layanan UMKM tidak cukup dibaca
sebagai masalah aplikasi, tetapi sebagai ketidakselarasan antara kapabilitas
akuisisi nasabah, penilaian kelayakan, tata kelola data, dan kemitraan digital.
Oleh karena itu, penelitian ini menggunakan capability map dan value stream
untuk mengidentifikasi kapabilitas prioritas sebelum merumuskan rancangan
arsitektur data, aplikasi, dan teknologi.
\end{quote}

\subsection{Dari Business Architecture ke Analisis}

Dalam bagian analisis, capability map dapat ditempatkan sebagai gambar atau
tabel. Value stream dapat digunakan untuk menjelaskan alur masalah. Yang
penting, artefak tidak berdiri sendiri. Setelah gambar atau tabel, selalu ada
paragraf yang menjelaskan temuan: kapabilitas mana yang lemah, pain point apa
yang muncul, dampaknya pada strategi, dan implikasinya bagi rancangan BDAT.

\section{Diskusi Kritis, Trade-off, dan Perangkap Umum}

\subsection{Trade-off dalam Business Architecture}

Business Architecture memaksa penulis memilih tingkat abstraksi. Jika terlalu
tinggi, capability map menjadi daftar kata besar yang sulit dipakai. Jika
terlalu rinci, peta berubah menjadi prosedur operasional dan kehilangan
fungsi strategis. Value stream juga memiliki trade-off serupa. Jika terlalu
pendek, pain point tidak terlihat. Jika terlalu panjang, pembaca kehilangan
alur utama.

\begin{table}[htbp]
\centering
\scriptsize
\caption{Trade-off utama dalam Business Architecture.}
\label{tab:tradeoff-bab4}
\begin{tabularx}{\textwidth}{@{}p{0.23\textwidth} X X X@{}}
\hline
\textbf{Keputusan} & \textbf{Alternatif} & \textbf{Manfaat} & \textbf{Risiko} \\
\hline
Cakupan capability map & Seluruh organisasi atau satu domain prioritas &
Cakupan luas memberi konteks; cakupan sempit lebih tajam & Terlalu luas
membuat makalah melebar. \\
Level detail kapabilitas & Kapabilitas besar atau subkapabilitas & Detail
membantu prioritas & Terlalu rinci membuat peta sulit dibaca. \\
Value stream & Alur end-to-end atau satu potongan kritis & End-to-end
menunjukkan konteks utuh & Terlalu panjang menyulitkan analisis. \\
Sumber bukti & Data internal, dokumen publik, atau asumsi eksplisit & Bukti
kuat memperkuat klaim & Data sensitif perlu izin dan anonimisasi. \\
Literature Review v1 & Ringkas terarah atau luas eksploratif & Ringkas lebih
mudah disintesis & Terlalu sempit dapat melewatkan konsep penting. \\
\hline
\end{tabularx}
\end{table}

\subsection{Perangkap Umum}

Perangkap pertama adalah menggambar capability map seperti bagan organisasi.
Kapabilitas bukan unit kerja. ``Divisi pemasaran'' adalah struktur organisasi;
``akuisisi nasabah digital'' adalah kapabilitas. Perangkap kedua adalah
menilai semua kapabilitas sebagai prioritas tinggi. Jika semua penting, maka
tidak ada strategi. Perangkap ketiga adalah membuat value stream sebagai
daftar aktivitas administratif tanpa menunjukkan nilai, aktor, pain point, dan
risiko.

Perangkap keempat adalah menulis Literature Review v1 sebagai ringkasan
artikel satu per satu. Tinjauan pustaka harus mensintesis. Artinya, penulis
perlu menunjukkan hubungan antar-konsep: bagaimana Enterprise Architecture
menjelaskan keselarasan strategi, bagaimana Business Architecture memetakan
kapabilitas, dan bagaimana value stream menunjukkan titik masalah yang perlu
dianalisis lebih lanjut.

\section{Latihan Terapan Individual}

Latihan berikut dikerjakan sendiri oleh setiap mahasiswa dan langsung menjadi
bahan makalah publikasi.

\subsection*{Latihan 4.1: Capability Map Studi Kasus}

Susun capability map untuk organisasi studi kasus Anda. Gunakan tiga sampai
empat kelompok kapabilitas. Tandai setiap kapabilitas dengan salah satu status:
dipertahankan, diperbaiki, atau dibangun. Setelah peta selesai, tulis satu
paragraf yang menjelaskan mengapa dua kapabilitas dipilih sebagai prioritas.

\subsection*{Latihan 4.2: Value Stream Prioritas}

Pilih satu value stream yang paling relevan dengan RQ. Tuliskan tahap utama
dari awal sampai akhir. Untuk setiap tahap, catat aktor, data yang dibutuhkan,
aplikasi atau kanal pendukung, pain point, dan risiko. Pilih tiga pain point
yang paling kuat untuk dibahas dalam makalah.

\subsection*{Latihan 4.3: Matriks Literatur untuk Business Architecture}

Ambil minimal lima referensi dari bibliografi awal Bab 3. Buat matriks yang
memuat konsep utama, fungsi dalam makalah, hubungan dengan capability map,
hubungan dengan value stream, dan kalimat sintesis yang mungkin dipakai.
Tambahkan referensi baru hanya bila benar-benar diperlukan dan dapat
diverifikasi.

\subsection*{Latihan 4.4: Draft Literature Review v1}

Tulis Literature Review v1 sepanjang 600 sampai 900 kata. Draf minimal memuat:
Enterprise Architecture sebagai konteks, Business Architecture sebagai alat
analisis, konsep kapabilitas atau value stream, dan satu paragraf sintesis
yang menjelaskan celah analisis pada studi kasus Anda. Draf ini belum final,
tetapi harus sudah memiliki sitasi yang benar.

\begin{table}[htbp]
\centering
\scriptsize
\caption{Rubrik mini latihan Bab 4.}
\label{tab:rubrik-bab4}
\begin{tabularx}{\textwidth}{@{}p{0.24\textwidth} X X@{}}
\hline
\textbf{Komponen} & \textbf{Indikator baik} & \textbf{Perlu diperbaiki bila} \\
\hline
Capability map & Kapabilitas relevan dengan RQ, dikelompokkan jelas, dan
memiliki prioritas & Peta menyerupai bagan organisasi atau terlalu umum. \\
Value stream & Alur nilai menunjukkan aktor, data, pain point, dan risiko &
Hanya menjadi daftar aktivitas tanpa analisis. \\
Prioritas & Alasan prioritas didukung bukti dan dampak & Semua kapabilitas
diberi prioritas sama. \\
Literature Review v1 & Sumber disintesis dan terhubung dengan artefak &
Hanya meringkas artikel satu per satu. \\
Narasi makalah & Artefak dijelaskan sebagai bukti analisis & Gambar atau tabel
tidak dijelaskan dalam teks. \\
\hline
\end{tabularx}
\end{table}

\subsection*{Refleksi Singkat}

Jawab tiga pertanyaan berikut:

\begin{enumerate}
  \item Kapabilitas apa yang paling menentukan keberhasilan topik makalah Anda?
  \item Pain point mana yang paling kuat sebagai bukti masalah penelitian?
  \item Konsep pustaka apa yang paling membantu menjelaskan temuan artefak?
\end{enumerate}

\section{Ringkasan, Refleksi, dan Jembatan}

\subsection*{Poin Kunci}

\begin{itemize}
  \item Business Architecture menerjemahkan visi arsitektur menjadi peta
        kemampuan bisnis yang dapat dianalisis.
  \item Capability map menunjukkan kapabilitas yang perlu dipertahankan,
        diperbaiki, atau dibangun.
  \item Value stream menunjukkan alur penciptaan nilai dan titik masalah yang
        memengaruhi pengguna, proses, data, aplikasi, dan risiko.
  \item Operating model membantu menentukan apa yang perlu distandardisasi dan
        apa yang perlu diintegrasikan.
  \item Literature Review v1 harus mensintesis konsep yang dipakai dalam
        capability map dan value stream.
  \item Artefak Business Architecture hanya berguna bila dijelaskan sebagai
        bagian dari argumen makalah.
\end{itemize}

\subsection*{Menjawab Pertanyaan Pemandu}

Kapabilitas bisnis yang paling penting adalah kapabilitas yang paling langsung
menghubungkan visi arsitektur dengan nilai bagi pemangku kepentingan.
Capability map membantu membedakan masalah inti dari gejala karena peta ini
menunjukkan kemampuan organisasi, bukan sekadar aplikasi atau unit kerja.
Value stream yang paling relevan adalah alur yang berhubungan langsung dengan
RQ dan menampilkan pain point yang dapat dianalisis. Bukti yang diperlukan
meliputi dokumen organisasi, observasi, wawancara, regulasi, dan literatur.
Hasil Business Architecture masuk ke makalah sebagai artefak analisis domain
bisnis dan sebagai dasar Literature Review v1.

\subsection*{Uji Mandiri}

\begin{enumerate}
  \item Jelaskan perbedaan antara kapabilitas bisnis dan proses bisnis.
  \item Mengapa capability map lebih stabil daripada bagan organisasi?
  \item Pilih satu kapabilitas pada studi kasus Anda. Data apa yang diperlukan
        untuk menilai kondisi awalnya?
  \item Apa perbedaan antara pain point dan akar masalah?
  \item Bagaimana value stream membantu menyiapkan analisis Data Architecture?
  \item Mengapa Literature Review v1 tidak cukup berisi ringkasan artikel?
\end{enumerate}

\subsection*{Jembatan ke Bab Berikutnya}

Bab berikutnya membahas Data Architecture. Hasil Bab 4 menjadi masukan
langsung untuk Bab 5. Capability map menunjukkan kapabilitas apa yang
membutuhkan data. Value stream menunjukkan titik aliran data, perpindahan data
antaraktor, dan risiko kualitas atau privasi. Dengan demikian, pembahasan data
pada Bab 5 tidak dimulai dari tabel database, melainkan dari kebutuhan bisnis
yang sudah terlihat dalam peta kapabilitas dan alur nilai.

\section{Bacaan Lanjutan dan Lampiran}

\subsection*{Bacaan Inti}

\begin{itemize}
  \item The TOGAF Standard untuk memahami posisi Business Architecture dalam
        pengembangan arsitektur enterprise \cite{togaf10}.
  \item Ross, Weill, dan Robertson untuk hubungan antara operating model,
        proses, data, dan eksekusi strategi \cite{ross2006enterprise}.
  \item Ross, Beath, dan Mocker untuk pembahasan kapabilitas organisasi digital
        dan cara merancang bisnis agar siap berubah \cite{ross2019designed}.
\end{itemize}

\subsection*{Bacaan untuk Memperdalam Topik}

\begin{itemize}
  \item Untuk topik platform dan ekosistem, gunakan literatur platform sebagai
        jembatan menuju Bab 11 \cite{parker2016platform,cusumano2019business}.
  \item Untuk topik integrasi layanan dan API, gunakan rujukan API sebagai
        pengantar menuju Bab 9 dan Bab 10 \cite{jacobson2011apis,medjaoui2018continuous}.
  \item Untuk topik modernisasi sistem, gunakan pembahasan arsitektur aplikasi
        dan microservices secara selektif \cite{newman2021building}.
  \item Untuk konteks Indonesia, pilih regulasi sesuai kasus: perlindungan data
        pribadi \cite{uupdp2022}, sistem elektronik \cite{pp712019}, Satu Data
        Indonesia \cite{satudata}, atau SNAP BI \cite{snapbi}.
\end{itemize}

\subsection*{Lampiran Kerja Bab 4}

Gunakan templat berikut sebagai catatan kerja:

\begin{description}
  \item[RQ dan visi arsitektur] Salin versi terbaru dari Bab 2 dan Bab 3.
  \item[Ruang lingkup bisnis] Tulis layanan, produk, unit, atau proses yang
        dianalisis.
  \item[Capability map] Kelompokkan kapabilitas dan beri status prioritas.
  \item[Value stream] Tuliskan tahap alur nilai, aktor, data, kanal, pain
        point, dan risiko.
  \item[Bukti] Catat sumber data atau asumsi untuk setiap klaim penting.
  \item[Lit Review v1] Tulis sintesis awal yang menghubungkan teori dengan
        artefak Business Architecture.
  \item[Implikasi Bab 5] Catat data apa yang perlu dipetakan pada Data
        Architecture.
\end{description}
\setlength{\emergencystretch}{0pt}
